\documentclass[11pt]{article}
\usepackage[margin=1in]{geometry}
\usepackage{amsmath,amssymb,amsthm,bm}
\usepackage{booktabs}
\usepackage{enumitem}
\usepackage[hidelinks]{hyperref}
\usepackage{xcolor}

\newtheorem*{example}{Example}
\newcommand{\code}[1]{\texttt{\small #1}}
\newcommand{\file}[1]{\par\noindent{\small\raggedright\color{gray} Implemented in \texttt{#1}.\par}}
\newcommand{\R}{\mathbb{R}}
\DeclareMathOperator*{\argmin}{arg\,min}

\title{The math behind the MHSeals boat tests}
\author{\texttt{mhseals\_testing}, \texttt{mhseals\_hardware}, \texttt{mhseals\_nav}}
\date{}

\begin{document}
\maketitle

\noindent This is a short course on every calculation the test suite and the
thruster chain make, in the order you meet them on the boat: from how a
command becomes four pulse widths, through fitting a model of the boat, to
scoring a navigation run. Each section ends with the file that implements it,
so you can read the math and the code side by side. Nothing here needs more
than vectors, matrices and a little calculus.

\tableofcontents
\bigskip

\section{Frames and notation}

\paragraph{Body frame.} Commands and velocities are in the boat's body frame,
\code{base\_link}, which follows ROS REP 103 (FLU):
$x$ forward (surge), $y$ to port (sway), and counter-clockwise rotation about
$z$ up (yaw). A command or velocity is the vector
\[
  \bm{u} = \begin{bmatrix} u_x & u_y & u_\psi \end{bmatrix}^\top ,\qquad
  \bm{v} = \begin{bmatrix} v_x & v_y & \omega_z \end{bmatrix}^\top .
\]
World frames (\code{map}, \code{odom}) are ENU: $x$ east, $y$ north, yaw
measured counter-clockwise from east.

\paragraph{Effort versus velocity.} There are two kinds of command, and
confusing them is the most common mistake:
\begin{itemize}[nosep]
  \item \emph{Effort} $\bm{e}\in[-1,1]^3$ is normalized thrust, what the
        thruster driver accepts. $e_x = 0.25$ means ``a quarter of full forward
        thrust'', not any particular speed.
  \item \emph{Velocity} $\bm{v}$ in m/s and rad/s is what Nav2 asks for. Turning
        velocity into effort needs a model of the boat (Section~\ref{sec:char})
        and a controller (Section~\ref{sec:controller}).
\end{itemize}
A test profile states which one its command topic carries
(\code{command.units}), and procedures that need one kind refuse to run with the
other.

\paragraph{Thrusters.} The four thrusters are indexed by position,
FL, FR, RR, RL, in that order everywhere. ``Physical output $k$'' is the
$k$-th PWM pin (or the $k$-th entry of the simulator's pulse array); the
\emph{channel map} says which output drives which position.
\file{mhseals\_hardware/thruster\_mixer.py, mhseals\_testing/profiles.py}

\section{The thrust mixer, from the boat's geometry}
\label{sec:mixer}

The mixer turns a body command (surge, sway, yaw) into one command per
thruster. It looks like a table of signs, but it comes straight out of where
the thrusters are and which way they push. This section derives it, explains
why it is only $\pm1$ for this boat, and what changes when the geometry is not
ideal.

\subsection{From thrusts to a body force}

Describe thruster $i$ by its position $(x_i, y_i)$ relative to the boat's
\emph{center of mass}, in the body frame ($x$ forward, $y$ to port), the
direction its positive thrust pushes the boat, an angle $\alpha_i$ from $+x$,
and its thrust $t_i$ (negative means it pushes the other way). Its force is
$t_i(\cos\alpha_i, \sin\alpha_i)$, and its turning moment is the 2-D cross
product of lever arm and force, $\tau_i = x_i F_{y,i} - y_i F_{x,i}$. Adding
the four thrusters,
\[
  \underbrace{\begin{bmatrix} F_x\\ F_y\\ \tau \end{bmatrix}}_{\bm{w}}
  = \underbrace{\begin{bmatrix}
      \cos\alpha_1 & \cdots & \cos\alpha_4\\
      \sin\alpha_1 & \cdots & \sin\alpha_4\\
      x_1\sin\alpha_1 - y_1\cos\alpha_1 & \cdots & x_4\sin\alpha_4 - y_4\cos\alpha_4
    \end{bmatrix}}_{B\ \in\ \R^{3\times4}}
  \underbrace{\begin{bmatrix} t_1\\ \vdots\\ t_4 \end{bmatrix}}_{\bm{t}} .
\]
Column $i$ of $B$ is what one unit of thrust from thruster $i$ does to the
boat. $B$ is the boat's \emph{allocation} matrix: kinematics, no tuning.

\subsection{From a body command to thrusts: the pseudo-inverse}

The mixer runs the other way: given a wanted $\bm{w}$, find $\bm{t}$. With four
thrusters and three axes, $B\bm{t} = \bm{w}$ has infinitely many solutions; any
$\bm{t}$ in the null space of $B$ (thrusts that cancel exactly) can be added
for free. The natural choice is the one that spends the least thrust,
\[
  \min_{\bm t}\ \sum_i t_i^2 \quad\text{subject to}\quad B\bm{t} = \bm{w}.
\]
The answer must lie in the row space of $B$, $\bm{t} = B^\top\bm{\lambda}$:
any component outside it lies in the null space, adds nothing to $\bm{w}$, and
only costs thrust. Substituting, $BB^\top\bm\lambda = \bm w$, so
\[
  \bm{t} = B^\top (BB^\top)^{-1}\bm{w} = B^{+}\bm{w},
\]
the Moore--Penrose pseudo-inverse. The transpose appears for the same reason
as in least squares (Section~\ref{sec:char}): the best solution is a
combination of the rows. $BB^\top$ is $3\times3$ and invertible whenever the
thrusters can produce every force and moment.

\subsection{Why this boat's mixer is only signs}

The boat is a symmetric ``X'': thrusters at $(\pm a, \pm b)$, each pushing
along a diagonal. Matching the code's sign convention,
\begin{center}
\begin{tabular}{lcccc}
  \toprule
  & FL & FR & RR & RL\\
  \midrule
  position & $(a, b)$ & $(a, -b)$ & $(-a, -b)$ & $(-a, b)$\\
  $\alpha_i$ & $135^\circ$ & $-135^\circ$ & $-45^\circ$ & $45^\circ$\\
  \bottomrule
\end{tabular}
\end{center}
Every $\cos\alpha_i$ and $\sin\alpha_i$ is $\pm\frac{1}{\sqrt2}$, and every
lever arm works out to $\pm\frac{a+b}{\sqrt2}$, so
\[
  B = \frac{1}{\sqrt2}\begin{bmatrix}
    -1 & -1 & +1 & +1\\
    +1 & -1 & -1 & +1\\
    (a{+}b) & -(a{+}b) & (a{+}b) & -(a{+}b)
  \end{bmatrix}.
\]
Its rows are orthogonal, so $BB^\top = \operatorname{diag}\!\left(2,\,2,\,2(a{+}b)^2\right)$
and $B^{+} = B^\top(BB^\top)^{-1}$ is just $B^\top$ with each column scaled.
Effort is normalized per axis anyway ($e=1$ means ``as much of this axis as one
saturated thruster allows''), so those scales drop out and the mixer is the
sign pattern of $B^\top$:
\[
  \bm{t} = M\bm{e},\qquad
  M = \begin{bmatrix}
    -1 & +1 & +1\\
    -1 & -1 & -1\\
    +1 & -1 & +1\\
    +1 & +1 & -1
  \end{bmatrix}
  \begin{matrix}\text{FL}\\ \text{FR}\\ \text{RR}\\ \text{RL}\end{matrix}
\]
Read a column to see how an axis is made: pure surge, $\bm e = (1,0,0)$, gives
$\bm t = (-1,-1,+1,+1)$; each of those thrusts points forward with a sideways
part, and the sideways parts cancel in pairs. Pure yaw drives diagonal pairs
against each other, so the forces cancel and only a moment remains. The
orthogonal rows are why $M^\top M = 4I$: asking for surge never adds sway or
yaw.

\paragraph{Desaturation.} A combined command can ask for more than one
thruster can give, e.g.\ $\bm{e} = (1,1,1)$ gives $t_\text{FR} = -3$.
Clipping that one thruster would change the \emph{direction} of the force.
Instead all four are scaled by the same factor,
\[
  \bm{t}' = \frac{\bm{t}}{\max\!\left(1, \max_i |t_i|\right)},
\]
which keeps the direction and gives up only magnitude.
\begin{example}
$\bm{e} = (1,1,1)$: $\bm{t} = (1,-3,1,1)$, peak $3$, so
$\bm{t}' = (\tfrac13, -1, \tfrac13, \tfrac13)$. The boat still moves forward,
to port and turns left, in the same proportions, at a third of the request.
\end{example}

\subsection{Other angles and rectangles}

\paragraph{Still mirror-symmetric.} Keep the X mirrored but let the thrusters
sit at any angle $\alpha$ from the long axis and at any rectangle $(\pm a,
\pm b)$. The surge row becomes $\pm\cos\alpha$, the sway row $\pm\sin\alpha$,
and every lever arm $\pm(a\sin\alpha + b\cos\alpha)$. The rows stay
orthogonal, so the $\pm1$ mixer still decouples the axes exactly; only the
\emph{strength} of each axis changes: surge capacity scales with $\cos\alpha$,
sway with $\sin\alpha$, yaw with the lever arm. This boat is already such a
rectangle: about $a \approx 0.27$\,m fore and aft and $b \approx 0.34$\,m to
the sides (hull offset plus mount), a yaw lever arm of
$(a+b)/\sqrt2 \approx 0.43$\,m. Unequal axis strengths need no new mixer:
the characterization (Section~\ref{sec:char}) measures them as the diagonal of
the gain $G$, and the velocity controller divides them out.

\paragraph{Asymmetric.} Symmetry breaks when the angles differ between
thrusters, the thrusters are not mirrored, or, most often, the center of mass
is not at the geometric center (batteries, a camera mast). Positions are
measured from the center of mass, so moving it by $(x_c, y_c)$ changes every
lever arm to $(x_i - x_c)\sin\alpha_i - (y_i - y_c)\cos\alpha_i$. The yaw row is
then no longer orthogonal to the surge and sway rows, and the $\pm1$ mixer
leaks: a pure sway command produces the moment
\[
  \tau = \sum_i t_i\,\big[(x_i - x_c)\sin\alpha_i - (y_i - y_c)\cos\alpha_i\big]
  \neq 0 \quad\text{for } \bm t = M(0,1,0)^\top,
\]
and the boat turns while sliding sideways. There are two ways to fix it:
\begin{enumerate}[nosep]
  \item \emph{From geometry:} measure $(x_i, y_i, \alpha_i)$ about the center of
        mass, build $B$, and use $M = B^{+}$, each column scaled so a full
        command just saturates the busiest thruster. Thrusters of unequal
        strength (or reverse thrust weaker than forward) call for a weighted
        pseudo-inverse: minimizing $\sum_i w_i t_i^2$ gives
        $\bm t = W^{-1}B^\top(BW^{-1}B^\top)^{-1}\bm w$ with $W = \operatorname{diag}(w_i)$.
  \item \emph{From data:} the characterization's off-diagonal gains are exactly
        this leakage, and the controller's $G^{-1}$ cancels it. Folding the
        correction into the mixer gives $M' = M G^{-1}$ (clipped to
        $[-1,1]$), the ``recommended allocation'' the fit reports; in effect it
        measures $B^{+}$ from how the boat moves instead of from a ruler.
\end{enumerate}
In practice use both: geometry for a good starting mixer, the fit for what the
geometry missed. A quick check is in the axis trials: pure surge should show
no steady yaw rate, and a consistent one means the yaw part of the mixer is off.
\file{mhseals\_hardware/thruster\_mixer.py (mix\_thrusters), mhseals\_testing/analysis/characterization.py}

\section{Effort to pulse width: the ESC model}
\label{sec:esc}

Each electronic speed controller (ESC) reads a pulse every 20\,ms (50\,Hz).
The pulse width $w$ in microseconds sets the thrust: $1500\,\mu$s is stopped,
longer is forward, shorter is reverse. The driver uses the linear map
\[
  w = \operatorname{clip}\!\left(1500 + 400\,t,\; 1100,\; 1900\right),
\]
so full effort $t=\pm1$ is $1900$ or $1100\,\mu$s. The pulse for thruster
position $p$ goes to physical output $c_p$ from the channel map.

The simulator inverts the same model. A pulse $w$ on an output becomes a motor
command
\[
  m = m_\text{max}\,\frac{w - 1500}{400},\qquad w\in[1100,1900],
\]
optionally with a deadband $|w-1500|\le d \Rightarrow m = 0$, which real ESCs
have. Because both sides share the model, a characterization done in simulation
exercises the same chain as on the boat.

\paragraph{Why republish.} The simulator treats a silent pulse topic as a lost
link and goes neutral after a timeout, as the real ESCs do without a signal. The
driver therefore repeats the last pulses at 20\,Hz even when they do not change.
\file{mhseals\_hardware/thruster\_mixer.py (effort\_to\_pulse), outputs.py; crane\_ml EscPwm.cs}

\section{Least-squares characterization}
\label{sec:char}

To command a velocity we need to know what velocity an effort produces. The
model is affine:
\[
  \bm{v} = \bm{b} + G\bm{e},
\]
where $G\in\R^{3\times3}$ is the \emph{gain} (m/s per unit effort on the
diagonal, cross-coupling off it) and $\bm{b}$ the \emph{bias}: the velocity at
zero effort, from current or wind.

\paragraph{Collecting samples.} The Characterize procedure holds a schedule of
efforts (each axis both ways at two levels, then combined directions), with
neutral between them. Every odometry sample $\bm{v}_k$ is paired with the effort
$\bm{e}_k$ held when it was measured, skipping the first second of each hold so
the fit sees the response and not the switching transient.

\paragraph{The fit.} Stack the samples as rows. For one output axis, say $v_x$,
\[
  \underbrace{\begin{bmatrix} v_{x,1}\\ \vdots\\ v_{x,n}\end{bmatrix}}_{\bm{y}}
  =
  \underbrace{\begin{bmatrix} 1 & \bm{e}_1^\top\\ \vdots & \vdots\\ 1 & \bm{e}_n^\top\end{bmatrix}}_{X}
  \underbrace{\begin{bmatrix} b_x\\ G_{x,\cdot}^\top\end{bmatrix}}_{\bm\theta}
  + \text{noise}.
\]
Least squares chooses $\bm\theta$ minimizing $\|\bm{y} - X\bm\theta\|^2$. Setting
the gradient to zero gives the \emph{normal equations}
\[
  X^\top X\,\bm\theta = X^\top \bm{y}.
\]
The $4\times4$ matrix $X^\top X$ is the same for all three axes, so one solve with
three right-hand sides gives $\bm{b}$ and every row of $G$.

\paragraph{When the fit is not trustworthy.} $X^\top X$ is invertible only if the
samples excite every column: some neutral, some surge, some sway, some yaw. The
code checks its rank and refuses a bag that only drove forward, rather than
returning a confident-looking but meaningless $G$. A tiny ridge term
$10^{-8} I$ guards against round-off.

\begin{example}
If $v_x = 0.6\,e_x$, $v_y = 0.4\,e_y$, $\omega_z = 0.3\,e_\psi$ exactly, the fit
returns $G = \operatorname{diag}(0.6, 0.4, 0.3)$ and $\bm b = \bm 0$. In the
simulator, $e_x = 0.2$ produced about $0.9$\,m/s, so $G_{xx}\approx 4.5$.
\end{example}

\paragraph{Limits.} The largest measured speed per axis gives the velocity
limits; the largest change of velocity between consecutive samples divided by
their time step gives acceleration and deceleration limits. These become the
Nav2 velocity-smoother limits, so Nav2 never asks for more than the boat showed
it can do.
\file{mhseals\_testing/analysis/characterization.py, procedures/characterize.py}

\section{The calibrated velocity controller}
\label{sec:controller}

Given a velocity request $\bm v^*$, invert the model to get the effort that
\emph{should} produce it (feed-forward), then correct what the model got wrong
with feedback on the measured velocity $\bm v$:
\[
  \bm e = \underbrace{G^{-1}\left(\operatorname{clip}(\bm v^*) - \bm b\right)}_{\text{feed-forward}}
        + K_p\,\bm\epsilon + K_i\,\bm z,\qquad \bm\epsilon = \bm v^* - \bm v,
\]
with $\bm e$ finally clipped to $[-1,1]$. $\operatorname{clip}$ limits the
request to the measured velocity limits.

\paragraph{Bounded integral (anti-windup).} The integral state
\[
  z_i \leftarrow \operatorname{clip}\!\left(z_i + \epsilon_i\,\Delta t,\; -z_\text{max},\; z_\text{max}\right)
\]
removes steady errors such as a constant current the bias did not capture. The
clip matters: without it, a long period of saturated thrust (a request the boat
cannot meet) keeps accumulating error, and the boat overshoots badly once it
finally can respond. That is called \emph{windup}. The integral also resets when
commands stop.

\paragraph{Why feed-forward first.} With a good model the feed-forward term does
almost all of the work, and the feedback gains can stay small and stable. If
the model is poor, the controller-steps procedure (Section~\ref{sec:step}) shows
it as large steady-state error or overshoot.
\file{mhseals\_hardware/velocity\_controller.py}

\section{Axis trials: is the boat moving the right way?}
\label{sec:trial}

An axis trial holds neutral for a baseline window, then a command on one axis,
then neutral again. Let $c$ be the velocity component for the axis
($v_x$, $v_y$ or $\omega_z$), $B$ the baseline samples and $D$ the driven samples.
\[
  \bar c_B = \frac1{|B|}\sum_{k\in B} c_k,\qquad
  \bar c_D = \frac1{|D|}\sum_{k\in D} c_k,\qquad
  r = \bar c_D - \bar c_B .
\]
The \emph{response} $r$ subtracts the drift already present before the command,
so a current pushing the boat forward does not make a reverse command look
weak. The verdict compares its sign with the commanded direction $s=\pm1$:
\[
  \text{verdict} = \begin{cases}
    \text{LOW RESPONSE} & |r| < 0.03\\
    \text{SIGN MATCH} & r\,s > 0\\
    \text{SIGN MISMATCH} & \text{otherwise.}
  \end{cases}
\]
\paragraph{When not to trust it.} The trial is INCONCLUSIVE with fewer than five
driven or three baseline samples, with a gap between samples longer than the
staleness limit, or if the odometry frames change mid-trial. A sample stamped at
the exact instant the command starts is not counted as driven: it cannot reflect
the command yet. SIGN MATCH says the wiring and signs are right; it is not a
calibration (that is Section~\ref{sec:char}).

\section{TF chain validation as a graph search}
\label{sec:tf}

ROS transforms form a tree: each frame has exactly one parent. Localization is
healthy only if the chains the stack needs, \code{map}$\to$\code{odom} and
\code{odom}$\to$\code{base\_link}, exist and are fresh.

Treat every recently received transform as a directed edge
parent$\to$child; static transforms never go stale. For each required chain
$(s, t)$, a depth-first search from $s$ looks for $t$. Three things can be wrong:
\begin{enumerate}[nosep]
  \item \emph{No path}: a node stopped publishing (MISSING).
  \item \emph{Two parents}: two nodes publish the same child, for example both
        the EKF and the simulator publishing \code{odom}$\to$\code{base\_link}.
        Edges are keyed by publisher identity, so this shows up as a frame with
        two parents on the path (WARN).
  \item \emph{A cycle}: a search that revisits a frame on its own path (WARN).
\end{enumerate}
Each search visits every edge at most once, so the check is linear in the
number of transforms.
\file{mhseals\_testing/monitor.py (\_tf\_rows)}

\section{Step-response metrics}
\label{sec:step}

The controller-steps procedure asks for a velocity step of size $r$ at time $0$
and records $y(t)$. Flip the sign for negative steps so that $r>0$.
\begin{description}[nosep]
  \item[Rise time] $t_{90} - t_{10}$, where $t_{p}$ is the first time $y\ge p\,r$.
  \item[Overshoot] $\max\!\left(0, \frac{\max_t y(t) - r}{r}\right)$.
  \item[Settling time] the first time after which $|y - r|$ stays within
        $\max(0.05\,r,\ 0.02)$. The floor keeps small steps from demanding
        precision below sensor noise.
  \item[Steady-state error] $r$ minus the mean of the last 20\% of the step;
        positive means the boat falls short.
  \item[IAE] the integral of absolute error, by the trapezoid rule:
        $\sum_k \tfrac12\left(|r-y_k| + |r-y_{k+1}|\right)(t_{k+1}-t_k)$.
\end{description}

\paragraph{Checking against theory.} For a first-order response
$y = r(1-e^{-t/\tau})$: no overshoot, rise time $\tau\ln 9$, and IAE exactly
$r\tau$. For an underdamped second-order system with damping ratio $\zeta$, the
overshoot is $e^{-\pi\zeta/\sqrt{1-\zeta^2}}$ and the 5\% settling time is about
$3/(\zeta\omega_n)$. The unit tests generate these curves and require the
metrics to match.
\file{mhseals\_testing/analysis/step\_response.py}

\section{Navigation: MPPI, the costmap, and how runs are measured}
\label{sec:nav}

\subsection{The MPPI controller}

Nav2 follows its planned path with MPPI (model predictive path integral). Every
control cycle it samples $K$ sequences of commands $(v_x, v_y, \omega)$ over a
horizon of $N$ steps of $\Delta t$, rolls each through a kinematic model of the
boat (Omni: the X thrusters can sway), scores each rollout with a sum of critics
$S_k$ (distance from the path, heading along it, nearness to obstacles, progress to
the goal, turning), and blends them with weights
\[
  w_k = \frac{e^{-(S_k - S_{\min})/\lambda}}{\sum_j e^{-(S_j - S_{\min})/\lambda}},
  \qquad \bm u = \sum_k w_k\,\bm u_k .
\]
The temperature $\lambda$ sets how much it averages: a high $\lambda$ blends many
samples, a low one commits to the best few. A floating boat's odometry is noisy
(waves) and its response lags, and averaging over many samples then washes out
the corrections; $\lambda = 0.1$ commits.

\paragraph{Timing.} After each cycle MPPI shifts its best sequence one step,
assuming a cycle lasts $\Delta t$. So $\Delta t$ must equal the control period:
$\Delta t = 1/f$. At $f = 20$\,Hz on a loaded computer the loop slipped to 11--17\,Hz,
the shifted plan fell out of step with the boat, and it drove in circles; at
$f = 10$\,Hz with $\Delta t = 0.1$\,s and $N=28$ (the same 2.8\,s horizon) it holds
its rate and drives straight.

\paragraph{Path horizon.} MPPI scores rollouts against the part of the path
within the prune distance $D$. A rollout runs up to $v_{\max} N\Delta t$; if that
exceeds $D$ the rollouts run off the end of the path, where only weak critics
apply, and near the end of a path the boat circles instead of arriving. Hence
$D \ge 1.2\, v_{\max} N \Delta t$.

\subsection{The costmap over open water}

A LiDAR over water returns almost nothing: the surface does not reflect the beam
back and there is nothing behind it. An obstacle grid clears a cell only when a
ray passes through it to a return, so a mark made once (a buoy, a gust-tilted
ring of points) would never be cleared. The obstacle layer therefore uses voxels
that decay: each mark lives a few seconds ($3$\,s in the local map, $15$\,s in the
global one, which must remember a buoy behind the boat) unless seen again. And a
VLP-16's rings are $2^\circ$ apart: a buoy $0.15$\,m above the water, $2$\,m away,
fits between them. The camera keeps such buoys in view, so the tracker's objects
are marked as small rings of points as a second source.

\paragraph{Inflation and the footprint.} The hull is a $1.06\times0.90$\,m
rectangle, padded by $0.1$\,m for tracking error at speed and for the scatter in
where a buoy is seen. Its padded corners lie
$r_c = \sqrt{0.63^2 + 0.55^2} \approx 0.84$\,m from the centre. MPPI scores
most trajectory points by the cost under the centre alone, and checks the whole
rectangle only where that cost says a corner could reach something, which is
the cost found at $r_c$. If the inflation ends short of $r_c$ that test can never
fire, and a corner can swing across a buoy the centre keeps clear of.

The inflation reaches further, to $1.5$\,m, and falls off gently, as
$252\,e^{-3(d - r_i)}$ beyond the inscribed radius $r_i = 0.55$\,m. Had it ended
at half a gate, the middle of any gap between two buoys a little wider than a
gate would cost nothing, and the planner, going from one task to the next, would
take the shortest line through the buoy field, sliding past a buoy with only the
footprint check's few centimetres to spare. Where a buoy is seen moves by
$0.1$ to $0.3$\,m while the boat turns, which is more than that. With the wider
inflation the gap costs something, and open water around the field is cheaper.
A gate's middle still has to be passable: in a 6\,ft ($1.83$\,m) gate it is
$0.91$\,m from each buoy and costs $252\,e^{-3(0.91-0.55)} \approx 84$, well
short of the inscribed cost, so the boat goes through when its goal lies beyond.

\subsection{Measuring a run}

\paragraph{Cross-track error.} The route is a polyline through the waypoints,
starting where the boat was when the goal was sent. For a segment from $\bm a$
to $\bm b$ and a position $\bm p$, project onto the segment and clamp to its ends:
\[
  \lambda = \operatorname{clip}\!\left(\frac{(\bm p-\bm a)\cdot(\bm b-\bm a)}{\|\bm b-\bm a\|^2},\,0,\,1\right),\qquad
  d = \left\|\bm p - (\bm a + \lambda(\bm b-\bm a))\right\|.
\]
The cross-track error of $\bm p$ is the smallest $d$ over all segments. Clamping
matters at corners: without it the distance would be to the infinite line, which
can be small far beyond the end of a leg. The procedure reports the mean, RMS and
maximum, the closest approach to each waypoint, and the travelled path length
against the route length.

\paragraph{Station keeping.} With samples $\bm p_k$ at times $t_k$ around a
centre $\bm c$, the radii $\rho_k = \|\bm p_k - \bm c\|$ give the RMS and maximum
radius, and the fraction of time inside the tolerance (each sample weighted by its
interval). A slow steady push that the controller does not cancel can stay inside
the tolerance for a long time, so the drift is measured directly as the
least-squares slope of each coordinate:
\[
  (\alpha_x, v_x) = \argmin_{\alpha,v} \sum_k \left(x_k - \alpha - v\,(t_k - t_0)\right)^2,
\]
and likewise for $y$. The drift speed is $\sqrt{v_x^2+v_y^2}$ and its heading
$\operatorname{atan2}(v_y, v_x)$.
\file{mhseals\_testing/analysis/nav\_metrics.py}

\section{Task geometry}
\label{sec:tasks}

\paragraph{Gates.} Red and green markers are paired greedily by distance: the
closest pair first, each marker used once, pairs wider than the maximum gate width
ignored. The gate is the midpoint $\bm m = \tfrac12(\bm r + \bm g)$. Its heading
is normal to the line $\bm r\to\bm g$; of the two normals, the one pointing away
from the boat is kept:
\[
  \phi = \operatorname{atan2}(g_y - r_y, g_x - r_x) + \tfrac\pi2,\qquad
  \text{if } \cos(\phi - \operatorname{atan2}(m_y - p_y, m_x - p_x)) < 0:\ \phi \leftarrow \phi + \pi,
\]
so the boat passes straight through instead of along the line of markers.

\paragraph{Circling.} To circle an object at $\bm o$ with radius $R$ in $n$ steps,
start on the side facing the boat, $\beta = \operatorname{atan2}(p_y - o_y, p_x - o_x)$,
and place waypoints at angles $\beta + s\,2\pi k/n$ for $k = 1\ldots n$, with
$s=+1$ counter-clockwise and $-1$ clockwise. Each waypoint faces along the
circle, angle $+\,s\,\pi/2$, so Nav2 drives around the object rather than at it.

\paragraph{Relative routes.} A route given as (forward, left) offsets
$(f, l)$ from a pose $(x, y, \psi)$ maps to the world by the rotation
$\left(x + f\cos\psi - l\sin\psi,\ y + f\sin\psi + l\cos\psi\right)$.
\file{mhseals\_nav/tasks/geometry.py}

\section{Mission scoring}
\label{sec:mission}

A task tree reports what it does as events on \code{/task/events}. A rubric
lists the events a good run must contain, with minimum counts $n_e$, the points
$p_e$ each event is worth, how many of an event count, $c_e$ (one unless the rubric
says more, so reporting the same thing twice earns nothing), and a time limit $T$. With $k_e$ the number of times event $e$ occurred:
\[
  \text{score} = \sum_e p_e \min(k_e, c_e),\qquad
  \text{possible} = \sum_e p_e\,c_e .
\]
The run passes when the task finished with SUCCESS, no required event is
missing ($k_e \ge n_e$ for all required $e$), and the elapsed time is at most
$T$. The report keeps the counts, missing events and points per event, so a
failed rehearsal says exactly what was not achieved.
\file{mhseals\_testing/analysis/mission.py, config/missions/}

\subsection{The referee}

The score above believes the task tree. In the simulator the referee checks what
the boat actually did, from the course the simulator placed (each object tagged
with its task) and the boat's true track.

\paragraph{Gates.} A gate is the segment between its two buoys. The track crossed
it between samples $\bm p$ and $\bm q$ when the segments intersect: the cross
products $c(\bm a,\bm b,\bm p)$ and $c(\bm a,\bm b,\bm q)$ differ in sign, and so do
$c(\bm p,\bm q,\bm a)$ and $c(\bm p,\bm q,\bm b)$, with
$c(\bm o,\bm a,\bm b) = (a_x-o_x)(b_y-o_y) - (a_y-o_y)(b_x-o_x)$. The sign of the
first says which way it crossed, so out and back through one gate shows as two
opposite crossings.

\paragraph{Contact.} With the hull a rectangle of length $L$ and beam $B$, a buoy
of radius $\rho$ at body-frame offset $(a, c)$ clears it by
$\sqrt{\max(|a| - L/2, 0)^2 + \max(|c| - B/2, 0)^2} - \rho$; a negative clearance
anywhere on the track is a contact.

\paragraph{Circling and docking.} The angle swept around the mark is the sum of
the wrapped changes of $\operatorname{atan2}(p_y - o_y, p_x - o_x)$ while within
$8$\,m of it; $300^\circ$ or more is a circle. Docked means the boat's centre
inside the right bay, facing into it within $30^\circ$, for the hold time.
\file{mhseals\_testing/analysis/referee.py}


\end{document}
